\subsection{拓扑的比较}

第 1 目的定理 2 表明，我们可以把连续映射取作拓扑结构的态射（《集合论》，第四章，§ 2，第 1 目）；从现在起，我们假定已作出了这样的态射选取。按照一般定义（《集合论》，第四章，§ 2，第 2 目），这使我们可以对给定集合 X 上的诸拓扑所成之集定义一个序：

定义 3. 给定同一集合 X 上的两个拓扑 \(\mathfrak{C}_1\)、\(\mathfrak{C}_2\)，我们说 \(\mathfrak{C}_1\) 比 \(\mathfrak{C}_2\) 细（而 \(\mathfrak{C}_2\) 比 \(\mathfrak{C}_1\) 粗），如果记 \(X_i\) 为赋以拓扑 \(\mathfrak{C}_i\)（\(i = 1, 2\)）的集合 X，恒等映射 \(X_1 \to X_2\) 连续。若此外还有 \(\mathfrak{C}_1 \neq \mathfrak{C}_2\)，则说 \(\mathfrak{C}_1\) 比 \(\mathfrak{C}_2\) 严格地细（而 \(\mathfrak{C}_2\) 比 \(\mathfrak{C}_1\) 严格地粗）。

两个拓扑，若其中一个比另一个细，则称它们是可比较的。

映射连续性的判别法（第 1 目，定义 1 与定理 1）给出下列命题：

命题 3. 给定两个拓扑 \(\mathfrak{G}_1\)、\(\mathfrak{G}_2\)（集合 \(X\) 上的拓扑），下列命题等价：

\begin{enumerate}
\def\labelenumi{\alph{enumi})}
\item
  \(\mathfrak{C}_1\) 比 \(\mathfrak{C}_2\) 细。
\item
  对每个 \(x \in \mathbf{X}\)，\(x\) 关于 \(\mathfrak{C}_2\) 的每个邻域都是 \(x\) 关于 \(\mathfrak{C}_1\) 的邻域。
\item
  对 X 的每个子集 A，A 在拓扑 \(C_{1}\) 中的闭包包含在 A 在拓扑 \(C_{2}\) 中的闭包之中。
\item
  在 \(\mathfrak{C}_{2}\) 中为闭集的 X 的每个子集在 \(\mathfrak{C}_{1}\) 中也是闭集。
\item
  在 \(\mathfrak{C}_{2}\) 中为开集的 X 的每个子集在 \(\mathfrak{C}_{1}\) 中也是开集。
\end{enumerate}

例。* 在由实数的序列 \(x = (x_n)\) 组成的 Hilbert 空间 H 中，这些序列满足

\[
\| x \| ^ {2} = \sum_ {n = 0} ^ {\infty} x _ {n} ^ {2} <   + \infty ,
\]

点 \(x_0\) 在 H 上的强拓扑中的邻域是包含开球 \(\|x - x_0\| < a\)（以 \(x_0\) 为中心的开球）的集合；\(x_0\) 在 H 上的弱拓扑中的邻域是包含由形如 \(\sup_{1 \leqslant i \leqslant n} |(x - x_0|a_i)| \leqslant 1\) 的关系所定义之集的集合，其中诸 \(a_i\) 是 H 的点，并且

\[
(x | y) = \sum_ {n = 0} ^ {\infty} x _ {n} y _ {n}
\]

（若 \(x = (x_n)\) 且 \(y = (y_n)\)）。现在若 \(\beta = \sup_{1 \leqslant i \leqslant n} \|a_i\|\)，则关系 \(\|x - x_0\| \leqslant 1/\beta\) 蕴含 \(|(x - x_0|a_i)| \leqslant \|x - x_0\| \cdot \|a_i\| \leqslant 1\)（对 \(1 \leqslant i \leqslant n\)）；因此 H 上的强拓扑比弱拓扑细。另一方面，对 H 的点的任意有限族 \((a_i)_{1 \leqslant i \leqslant n}\)，存在 H 中的点 \(x\)，使 \((x - x_0|a_i) = 0\)（对 \(1 \leqslant i \leqslant n\)）且使 \(\|x - x_0\|\) 任意大；这表明强拓扑比弱拓扑严格地细。*

注。1) 在集合 X 上所有拓扑组成的有序集中，仅以 \(\varnothing\) 与 X 为开集的拓扑是最粗的，而离散拓扑是最细的。

\begin{enumerate}
\def\labelenumi{\arabic{enumi})}
\setcounter{enumi}{1}
\item
  拓扑越细，开集、闭集与邻域就越多；拓扑越细，集合的闭包（相应地，内部）就越小（相应地，越大）；拓扑越细，稠密集就越少。
\item
  如果 \(f: \mathbf{X} \to \mathbf{X}'\) 是连续映射，那么当把 \(\mathbf{X}\) 的拓扑换成较细的拓扑、把 \(\mathbf{X}'\) 的拓扑换成较粗的拓扑时，它仍然连续（第 1 目，定理 2）。换言之，\(\mathbf{X}\) 的拓扑越细且 \(\mathbf{X}'\) 的拓扑越粗，由 \(\mathbf{X}\) 到 \(\mathbf{X}'\) 中的连续映射就越多。
\end{enumerate}

